\documentclass[11pt,letterpaper]{report}

\usepackage[T1]{fontenc}
\usepackage[utf8]{inputenc}
\usepackage[margin=1in]{geometry}
\usepackage{booktabs}
\usepackage{xltabular}
\usepackage[colorlinks=true,linkcolor=blue,urlcolor=blue]{hyperref}
\usepackage[expansion=false]{microtype}
\usepackage{makecell}
\usepackage{parskip}
\usepackage[explicit]{titlesec}

% Chapter heading style
\titleformat{\chapter}[hang]
  {\normalfont\LARGE\bfseries}
  {\thechapter.}{0.5em}{#1}
\titlespacing*{\chapter}{0pt}{-10pt}{20pt}

% Section heading
\titleformat{\section}[hang]
  {\normalfont\large\bfseries}
  {\thesection}{0.5em}{#1}

\begin{document}
\begin{titlepage}
\vspace*{2cm}
{\Huge\bfseries SVX Applications Catalog\par}
\vspace{0.5cm}
{\large SystemVerilog Extension Library --- Applications (\texttt{apps/})\par}
\vspace{1cm}
{\large Version 2.1.2~beta\par}
\vspace{0.3cm}
{\large September 2026\par}
\vspace{0.3cm}
{\large Mark Glasser\par}
\vspace{1cm}
{\small Copyright 2026 Mark Glasser.
Licensed under the Apache License, Version~2.0.\par}
\end{titlepage}

\tableofcontents
\newpage
\chapter{Sparse Memory}


The \texttt{mem} application provides a simulation-friendly model of a large
address space without requiring backing storage for every possible location.
Memory is allocated lazily: pages and blocks are created on demand as data
is written, so only those locations that actually receive writes consume
simulator memory.  A sparse representation makes it practical to model a
32-bit (or wider) address space in a verification environment without
pre-allocating four~gigabytes of simulator memory.

\section*{Address Decomposition}

Every address is decomposed into three fields according to parameterized
bit widths.  The upper \texttt{PAGE\_BITS} bits identify the \emph{page};
the next \texttt{BLOCK\_BITS} bits identify the \emph{block} within the
page; the remaining low-order bits identify the \emph{byte} within the
block.  This three-level hierarchy mirrors the structure of a real memory
subsystem and makes sparse allocation efficient even for large address
spaces.

\section*{Security Model}

The sparse memory supports read/write restrictions at three levels of
granularity: page, block, and individual word.  A restriction value of
\texttt{RESTRICT\_READ} prevents reads; \texttt{RESTRICT\_WRITE} prevents
writes; \texttt{RESTRICT\_READ\_WRITE} prevents both.  Restrictions are
consulted in order from coarsest (page) to finest (word); the first
matching restriction wins.  When a security violation occurs the offending
operation is aborted and an error code is recorded.

\section*{Error Reporting}

After every read or write operation the caller can inspect whether the
operation succeeded.  \texttt{last\_operation\_failed()} returns
\texttt{1} if an error occurred.  \texttt{get\_last\_error()} returns the
specific \texttt{error\_t} value, and \texttt{get\_last\_op\_string()}
returns a human-readable string that names the operation, the address, and
the reason for failure.  This is particularly useful for debug and
self-checking testbenches.

\section*{Bounded Memory}

The \texttt{mem\_bounded} class extends \texttt{mem} with configurable
lower and upper address bounds.  Any access outside those bounds is
rejected with an \texttt{ERROR\_ADDRESS\_BOUNDS} error.  An optional
bounds lock prevents the bounds from being changed once set, which is
useful for enforcing region ownership in a shared address space.

\section*{Implementation Notes}

Internally the sparse memory is structured as a \emph{map of pages}, where
each page is a \emph{map of blocks}, and each block is a \emph{vector of
bytes}.  The maps use SVX's \texttt{map} class (a balanced tree), so
lookup is $O(\log n)$ in the number of allocated pages or blocks.  Only the
\texttt{mem} and \texttt{mem\_bounded} classes are intended for direct
instantiation; \texttt{mem\_page} and \texttt{mem\_block} are internal
implementation classes.

\section{Parameters}

All four classes share the same parameter set.

\begin{xltabular}{\linewidth}{llX}
\toprule
\textbf{Parameter} & \textbf{Default} & \textbf{Description} \\
\midrule
\endfirsthead
\toprule
\textbf{Parameter} & \textbf{Default} & \textbf{Description} \\
\midrule
\endhead
\bottomrule
\endlastfoot
\texttt{ADDR\_BITS} & \texttt{32} & Total address space width in bits. \\
\texttt{PAGE\_BITS} & \texttt{16} & Width of the page index field. \\
\texttt{BLOCK\_BITS} & \texttt{8} & Width of the block index field within a page. \\
\texttt{WORD\_SIZE} & \texttt{4} & Word size in bytes; words must be naturally aligned. \\
\end{xltabular}

\section{Types}

Key typedefs (defined in \texttt{mem\_types.svh} and \texttt{mem\_base}):

\begin{xltabular}{\linewidth}{lX}
\toprule
\textbf{Type} & \textbf{Description} \\
\midrule
\endfirsthead
\toprule
\textbf{Type} & \textbf{Description} \\
\midrule
\endhead
\bottomrule
\endlastfoot
\texttt{addr\_t} & Parameterized address type: bit [ADDR\_BITS-1:0]. \\
\texttt{word\_t} & Word type: bit [WORD\_SIZE*8-1:0]. \\
\texttt{byte\_t} & 8-bit unsigned byte. \\
\texttt{pri\_t} & Priority type (uint32\_t), used by pri\_queue, not mem. \\
\texttt{restrict\_t} & Security restriction enum (see below). \\
\texttt{error\_t} & Error code enum (see below). \\
\texttt{operation\_t} & Last operation enum (see below). \\
\end{xltabular}

\noindent\textbf{\texttt{restrict\_t}}

\begin{xltabular}{\linewidth}{lX}
\toprule
\textbf{Value} & \textbf{Meaning} \\
\midrule
\endfirsthead
\toprule
\textbf{Value} & \textbf{Meaning} \\
\midrule
\endhead
\bottomrule
\endlastfoot
\texttt{RESTRICT\_NONE} & No restriction; reads and writes are allowed. \\
\texttt{RESTRICT\_READ} & Reads are blocked; writes are allowed. \\
\texttt{RESTRICT\_WRITE} & Writes are blocked; reads are allowed. \\
\texttt{RESTRICT\_READ\_WRITE} & Both reads and writes are blocked. \\
\end{xltabular}

\noindent\textbf{\texttt{error\_t}}

\begin{xltabular}{\linewidth}{lX}
\toprule
\textbf{Value} & \textbf{Meaning} \\
\midrule
\endfirsthead
\toprule
\textbf{Value} & \textbf{Meaning} \\
\midrule
\endhead
\bottomrule
\endlastfoot
\texttt{ERROR\_NONE} & No error; the last operation succeeded. \\
\texttt{ERROR\_PARAMETERS\_WRONG} & Class parameters are inconsistent (fatal). \\
\texttt{ERROR\_MISALIGNMENT} & Word access to a non-word-aligned address. \\
\texttt{ERROR\_PAGE\_SECURITY\_VIOLATION} & A page-level restriction blocked the access. \\
\texttt{ERROR\_BLOCK\_SECURITY\_VIOLATION} & A block-level restriction blocked the access. \\
\texttt{ERROR\_WORD\_SECURITY\_VIOLATION} & A word-level restriction blocked the access. \\
\texttt{ERROR\_ADDRESS\_BOUNDS} & Address is outside the configured bounds (mem\_bounded only). \\
\end{xltabular}

\noindent\textbf{\texttt{operation\_t}}

\begin{xltabular}{\linewidth}{lX}
\toprule
\textbf{Value} & \textbf{Meaning} \\
\midrule
\endfirsthead
\toprule
\textbf{Value} & \textbf{Meaning} \\
\midrule
\endhead
\bottomrule
\endlastfoot
\texttt{OP\_NONE} & No operation has been performed yet. \\
\texttt{OP\_READ} & The last operation was a word read. \\
\texttt{OP\_WRITE} & The last operation was a word write. \\
\texttt{OP\_READ\_BYTE} & The last operation was a byte read. \\
\texttt{OP\_WRITE\_BYTE} & The last operation was a byte write. \\
\end{xltabular}

\section{Class: \texttt{mem}}

The top-level user-facing class.  Declare and instantiate it directly:

\texttt{mem\#(ADDR\_BITS, PAGE\_BITS, BLOCK\_BITS, WORD\_SIZE) m = new();}

\begin{xltabular}{\linewidth}{lllX}
\toprule
\textbf{Method} & \textbf{Returns} & \textbf{Arguments} & \textbf{Description} \\
\midrule
\endfirsthead
\toprule
\textbf{Method} & \textbf{Returns} & \textbf{Arguments} & \textbf{Description} \\
\midrule
\endhead
\bottomrule
\endlastfoot
\texttt{write} & void & \makecell[tl]{\texttt{addr\_t addr} \\ \texttt{word\_t data}} & Write a word to addr. Addr must be word-aligned. \\
\texttt{read} & \texttt{word\_t} & \texttt{addr\_t addr} & Read and return the word at addr. Returns 0 if the location was never written. \\
\texttt{write\_byte} & void & \makecell[tl]{\texttt{addr\_t addr} \\ \texttt{byte\_t data}} & Write a single byte to addr. \\
\texttt{read\_byte} & \texttt{byte\_t} & \texttt{addr\_t addr} & Read and return the byte at addr. Returns 0 if never written. \\
\texttt{set\_page\_restriction} & void & \makecell[tl]{\texttt{addr\_t addr} \\ \texttt{restrict\_t r}} & Apply restriction r to the entire page that contains addr. \\
\texttt{get\_page\_restriction} & \texttt{restrict\_t} & \texttt{addr\_t addr} & Return the restriction on the page containing addr. \\
\texttt{clear\_page\_restriction} & void & \texttt{addr\_t addr} & Remove any page-level restriction for the page containing addr. \\
\texttt{set\_block\_restriction} & void & \makecell[tl]{\texttt{addr\_t addr} \\ \texttt{restrict\_t r}} & Apply restriction r to the block containing addr. \\
\texttt{clear\_block\_restriction} & void & \texttt{addr\_t addr} & Remove any block-level restriction for the block containing addr. \\
\texttt{set\_word\_restriction} & void & \makecell[tl]{\texttt{addr\_t addr} \\ \texttt{restrict\_t r}} & Apply restriction r to the word at addr. \\
\texttt{clear\_word\_restriction} & void & \texttt{addr\_t addr} & Remove any word-level restriction at addr. \\
\texttt{is\_readable} & \texttt{bit} & \texttt{addr\_t addr} & Return 1 if addr may be read (no active read restriction). \\
\texttt{is\_writable} & \texttt{bit} & \texttt{addr\_t addr} & Return 1 if addr may be written (no active write restriction). \\
\texttt{get\_addr\_restriction} & \texttt{restrict\_t} & \texttt{addr\_t addr} & Return the effective restriction at addr (coarsest level that applies). \\
\texttt{last\_operation\_failed} & \texttt{bit} & --- & Return 1 if the most recent operation encountered an error. \\
\texttt{get\_last\_error} & \texttt{error\_t} & --- & Return the error\_t code from the most recent operation. \\
\texttt{get\_last\_op\_string} & \texttt{string} & --- & Return a human-readable string describing the last operation and any error. \\
\texttt{dump\_security} & void & \texttt{addr\_t addr=0} & Print the security map for all pages (and blocks if allocated). \\
\texttt{dump} & void & \texttt{addr\_t addr=0} & Print the contents of all allocated pages and blocks. \\
\end{xltabular}
\section{Class: \texttt{mem\_bounded}}

\texttt{mem\_bounded} extends \texttt{mem} with lower and upper address bounds. Any access outside the configured range is rejected with \texttt{ERROR\_ADDRESS\_BOUNDS}. All methods of \texttt{mem} are available; the following are added or overridden.

\begin{xltabular}{\linewidth}{lllX}
\toprule
\textbf{Method} & \textbf{Returns} & \textbf{Arguments} & \textbf{Description} \\
\midrule
\endfirsthead
\toprule
\textbf{Method} & \textbf{Returns} & \textbf{Arguments} & \textbf{Description} \\
\midrule
\endhead
\bottomrule
\endlastfoot
\texttt{set\_bounds} & void & \makecell[tl]{\texttt{addr\_t lower} \\ \texttt{addr\_t upper}} & Set both the lower and upper bounds in one call. \\
\texttt{set\_lower\_bound} & void & \texttt{addr\_t lower} & Set the lower bound only. \\
\texttt{set\_upper\_bound} & void & \texttt{addr\_t upper} & Set the upper bound only. \\
\texttt{get\_lower\_bound} & \texttt{addr\_t} & --- & Return the current lower bound. \\
\texttt{get\_upper\_bound} & \texttt{addr\_t} & --- & Return the current upper bound. \\
\texttt{set\_bounds\_lock} & void & \texttt{bit lock} & Lock the address bounds when lock=1; further calls to set\_bounds* are silently ignored. \\
\texttt{get\_bounds\_lock} & \texttt{bit} & --- & Return 1 if the bounds are currently locked. \\
\end{xltabular}
\chapter{Memory Map}


The \texttt{mem\_map} application models a hierarchical address space of the
kind found in hardware block specifications: registers, bit fields,
sub-regions, memories, and overlapping views, all organized into a tree
whose nodes carry precise offset and size information.

\section*{Overview}

Each node in the hierarchy is a \emph{memory space}, represented by the
abstract base class \texttt{mem\_space}.  Concrete subclasses capture the
five distinct space types recognized in hardware design: \texttt{MEMORY},
\texttt{REGISTER}, \texttt{FIELD}, \texttt{REGION}, and \texttt{VIEW}.
Every space has a name, a parent, an offset (relative to the parent), and a
size.  The hierarchy is built by constructing child spaces and passing the
parent as a constructor argument; the child registers itself with its parent
automatically.

\section*{Space Types and Hierarchy Rules}

The five space types form a strict containment hierarchy.  Not every type
may be nested inside every other type:

\begin{itemize}
  \item A \textbf{FIELD} is a bit field within a register.  Fields are
        leaves: they may not contain children.
  \item A \textbf{REGISTER} models a memory-mapped register.  Its only
        legal children are FIELDs.
  \item A \textbf{MEMORY} represents a contiguous block of memory.
        Memories are leaves: they may not contain children.
  \item A \textbf{REGION} is a container for other spaces.  It may
        contain REGISTERs, MEMORies, nested REGIONs, and VIEWs, but
        not FIELDs directly.  Children of a region may not overlap
        each other.
  \item A \textbf{VIEW} is like a REGION, but sibling VIEWs \emph{are}
        permitted to overlap each other.  This models hardware designs
        where the same physical addresses can be interpreted in different
        ways depending on a mode bit.
\end{itemize}

\section*{Address Calculation and Validation}

After the hierarchy is fully constructed, call
\texttt{calculate\_and\_check()} on the root to propagate absolute
addresses down the tree and to check for overlapping siblings.
\texttt{get\_addr()} and \texttt{get\_end\_addr()} then return the
absolute address range of any node.

\section*{Navigation}

\texttt{find\_addr(addr)} searches the hierarchy for the deepest node
whose range contains \texttt{addr} and returns it.
\texttt{find\_addr\_all(addr)} returns a list of \emph{all} nodes (at
every level) whose range contains \texttt{addr}, which is useful when
VIEWs overlap.

\section*{Implementation}

\texttt{mem\_space} extends SVX's \texttt{tree} class.  Children are
stored in an SVX \texttt{deque} and accessed through a typed iterator.
The overlap check in \texttt{check()} iterates over all pairs of
non-VIEW siblings and verifies that their address ranges do not
intersect.

\section{Space Types}

\begin{xltabular}{\linewidth}{llX}
\toprule
\textbf{Type} & \textbf{Leaf?} & \textbf{Allowed children} \\
\midrule
\endfirsthead
\toprule
\textbf{Type} & \textbf{Leaf?} & \textbf{Allowed children} \\
\midrule
\endhead
\bottomrule
\endlastfoot
\texttt{FIELD} & Yes & None \\
\texttt{REGISTER} & No & FIELD only \\
\texttt{MEMORY} & Yes & None \\
\texttt{REGION} & No & REGISTER, MEMORY, REGION, VIEW \\
\texttt{VIEW} & No & REGISTER, MEMORY, REGION, VIEW (may overlap sibling VIEWs) \\
\end{xltabular}

\section{Class: \texttt{mem\_space} (Abstract Base)}

All concrete space classes extend \texttt{mem\_space\#(ADDR\_SIZE=32)}. The methods below are defined on the base and inherited by all subtypes. Pure-virtual methods (\texttt{get\_type()}, \texttt{check\_child()}) are implemented by each concrete subclass.

\begin{xltabular}{\linewidth}{lllX}
\toprule
\textbf{Method} & \textbf{Returns} & \textbf{Arguments} & \textbf{Description} \\
\midrule
\endfirsthead
\toprule
\textbf{Method} & \textbf{Returns} & \textbf{Arguments} & \textbf{Description} \\
\midrule
\endhead
\bottomrule
\endlastfoot
\texttt{get\_offset} & \texttt{addr\_t} & --- & Return the offset of this space relative to its parent. \\
\texttt{get\_size} & \texttt{mem\_size\_t} & --- & Return the size of this space in address units. \\
\texttt{get\_addr} & \texttt{addr\_t} & --- & Return the absolute base address (valid after calculate\_and\_check()). \\
\texttt{get\_end\_addr} & \texttt{addr\_t} & --- & Return the absolute end address (base + size - 1). \\
\texttt{get\_type} & \texttt{mem\_space\_type\_t} & --- & Return the space type (FIELD, REGISTER, MEMORY, REGION, or VIEW). Pure virtual. \\
\texttt{get\_type\_name} & \texttt{string} & --- & Return the space type as a human-readable string. \\
\texttt{check\_child} & \texttt{bit} & \texttt{mem\_space\_t child} & Return 1 if child is a legal child of this space type. Pure virtual. \\
\texttt{insert\_space} & void & \texttt{mem\_space\_t child} & Insert child into this space; emits a warning and skips insertion on type violation. \\
\texttt{get\_error} & \texttt{bit} & --- & Return 1 if this space has been flagged as erroneous. \\
\texttt{clear\_error} & void & --- & Clear the stored error message. \\
\texttt{calculate\_and\_check} & \texttt{bit} & --- & Propagate absolute addresses and check for overlapping siblings. Call on the root after the hierarchy is fully built. \\
\texttt{find\_addr} & \texttt{mem\_space\_t} & \texttt{addr\_t addr} & Find and return the deepest space whose range contains addr. \\
\texttt{find\_addr\_all} & \texttt{list\_t} & \texttt{addr\_t addr} & Return a list of all spaces (at every level) whose range contains addr. \\
\texttt{dump} & void & --- & Print the hierarchy rooted at this node, with offsets, sizes, and absolute addresses. \\
\end{xltabular}
\section{Classes: \texttt{mem\_region} and \texttt{mem\_view}}

\texttt{mem\_region\#(ADDR\_SIZE=32)} and \texttt{mem\_view\#(ADDR\_SIZE=32)} both extend \texttt{mem\_space} and provide factory methods for building the hierarchy.  \texttt{mem\_view} additionally permits sibling instances to overlap in address space.  Constructor: \texttt{new(name, parent, offset, size)}.

\begin{xltabular}{\linewidth}{lllX}
\toprule
\textbf{Method} & \textbf{Returns} & \textbf{Arguments} & \textbf{Description} \\
\midrule
\endfirsthead
\toprule
\textbf{Method} & \textbf{Returns} & \textbf{Arguments} & \textbf{Description} \\
\midrule
\endhead
\bottomrule
\endlastfoot
\texttt{add\_register} & void & \makecell[tl]{\texttt{string name} \\ \texttt{addr\_t offset} \\ \texttt{mem\_size\_t size}} & Create and attach a mem\_register child at the given offset and size. \\
\texttt{add\_memory} & void & \makecell[tl]{\texttt{string name} \\ \texttt{addr\_t offset} \\ \texttt{mem\_size\_t size}} & Create and attach a mem\_memory child at the given offset and size. \\
\texttt{add\_region} & void & \makecell[tl]{\texttt{string name} \\ \texttt{addr\_t offset} \\ \texttt{mem\_size\_t size}} & Create and attach a nested mem\_region child. \\
\texttt{add\_view} & void & \makecell[tl]{\texttt{string name} \\ \texttt{addr\_t offset} \\ \texttt{mem\_size\_t size}} & Create and attach a mem\_view child. \\
\end{xltabular}
\section{Class: \texttt{mem\_register}}

\texttt{mem\_register\#(ADDR\_SIZE=32)} extends \texttt{mem\_space}. Its only legal children are FIELDs. Constructor: \texttt{new(name, parent, offset, size)}.

\begin{xltabular}{\linewidth}{lllX}
\toprule
\textbf{Method} & \textbf{Returns} & \textbf{Arguments} & \textbf{Description} \\
\midrule
\endfirsthead
\toprule
\textbf{Method} & \textbf{Returns} & \textbf{Arguments} & \textbf{Description} \\
\midrule
\endhead
\bottomrule
\endlastfoot
\texttt{add\_field} & void & \makecell[tl]{\texttt{string name} \\ \texttt{addr\_t offset} \\ \texttt{mem\_size\_t size}} & Create and attach a mem\_field child representing a bit field at the given bit offset and width. \\
\end{xltabular}
\section{Classes: \texttt{mem\_memory} and \texttt{mem\_field}}

Both are leaf nodes with no children and no additional methods beyond those inherited from \texttt{mem\_space}.  They exist primarily as type tags so that \texttt{check\_child()} can enforce the containment rules. Constructor for both: \texttt{new(name, parent, offset, size)}.

\chapter{Priority Queue}


The \texttt{pri\_queue} application is a generic priority queue that
supports two fundamental operations: \texttt{push()} and \texttt{pop()}.

\section*{Overview}

\texttt{push(pri, item)} inserts \texttt{item} into the queue at priority
level \texttt{pri}.  \texttt{pop()} removes and returns the item with the
highest priority.  When multiple items share the same priority they are
returned in the order they were inserted --- first-in, first-out within
each priority level.

\section*{Implementation}

The underlying data structure is a \emph{map of queues}: an SVX
\texttt{map} whose keys are priority values and whose values are SVX
\texttt{queue} instances holding the items at each priority.  Because SVX's
\texttt{map} is implemented as an ordered tree, the entry with the highest
key (priority) is retrieved in constant time via \texttt{last()}.  When the
last item at a given priority is popped, the corresponding map entry is
deleted automatically, keeping the map compact.

\section*{Type Parameters}

\texttt{pri\_queue} is doubly parameterized: \texttt{T} is the item type
(default \texttt{int}) and \texttt{P} is the traits class (default
\texttt{void\_traits}).  The priority type \texttt{pri\_t} is a
\texttt{uint32\_t} defined in \texttt{pri\_queue\_pkg}.  \texttt{pri\_queue}
extends SVX's \texttt{typed\_container\#(T,P)} and inherits its
\texttt{is\_empty()}, \texttt{size()}, and \texttt{clear()} methods.

\section{Class: \texttt{pri\_queue}}

Declared as \texttt{pri\_queue\#(type T=int, type P=void\_traits)} in \texttt{pri\_queue\_pkg}.  Constructor: \texttt{new()} --- no arguments required.

\begin{xltabular}{\linewidth}{lllX}
\toprule
\textbf{Method} & \textbf{Returns} & \textbf{Arguments} & \textbf{Description} \\
\midrule
\endfirsthead
\toprule
\textbf{Method} & \textbf{Returns} & \textbf{Arguments} & \textbf{Description} \\
\midrule
\endhead
\bottomrule
\endlastfoot
\texttt{push} & void & \makecell[tl]{\texttt{pri\_t pri} \\ \texttt{T item}} & Insert item into the queue at priority level pri. Higher values of pri represent higher priority. \\
\texttt{pop} & \texttt{T} & --- & Remove and return the item with the highest priority. When multiple items share the top priority, returns the one inserted earliest. Returns P::empty if the queue is empty. \\
\texttt{is\_empty} & \texttt{bit} & --- & Return 1 if the queue contains no items. \\
\texttt{size} & \texttt{size\_t} & --- & Return the total number of items across all priority levels. \\
\texttt{clear} & void & --- & Remove all items from the queue. \\
\end{xltabular}
\end{document}
